\chapter{Empirical Strategy}
\label{chap:methods}

\section{Estimands before models}

The phrase ``populism is associated with corruption'' can describe three different quantities. A
pooled estimate combines differences between countries and changes within countries. A between
estimate asks whether countries with more populist governments are more corrupt on average. A fixed-
effects estimate asks whether a country is more corrupt in years when its own government is more
populist than usual. Finally, a directional model asks whether past values of one variable predict
later values of the other after persistence is considered.

No single model answers all these questions. The analysis therefore fixes the estimand first and uses
model fit only as a diagnostic. In particular, a pooled regression may predict levels well because it
uses large cross-country differences, yet remain inappropriate for the within-country research
question. Conversely, a low within-
\Rtwo{} does not invalidate a fixed-effects coefficient when the dependent variable is highly
persistent and the purpose is estimation rather than prediction \citep{wooldridge2010}.

\section{Contemporaneous association}

\subsection{Two-way fixed effects}

The primary association model is

\begin{equation}
  C_{it} = \beta P_{it} + \mathbf{X}_{it}'\boldsymbol{\gamma}
  + \alpha_i + \lambda_t + \varepsilon_{it},
  \label{eq:twfe}
\end{equation}

where $C_{it}$ is composite corruption, $P_{it}$ is governing-party populism,
$\mathbf{X}_{it}$ contains time-varying controls, $\alpha_i$ are country fixed effects, and
$\lambda_t$ are year fixed effects. Country effects remove time-invariant history, geography, and
baseline institutions. Year effects remove shocks common to all countries. The coefficient $\beta$
is therefore identified by within-country deviations from long-run country and year means.

Three static specifications are fitted on one frozen complete-case sample:

\begin{itemize}
  \item M1: populism plus country and year effects;
  \item M2: M1 plus log GDP per capita;
  \item M3: M2 plus electoral democracy.
\end{itemize}

Freezing the M1--M3 sample ensures that coefficient movement reflects added controls rather than a
changing set of observations. Country-clustered standard errors allow arbitrary serial dependence
within countries. Two-way clustering and Driscoll--Kraay covariance estimates provide sensitivity to
common cross-sectional shocks \citep{driscoll1998}.

M4 adds lagged corruption:

\begin{equation}
  C_{it} = \rho C_{i,t-1} + \beta_{SR} P_{it}
  + \mathbf{X}_{it}'\boldsymbol{\gamma} + \alpha_i + \lambda_t + u_{it}.
  \label{eq:dynamic}
\end{equation}

This is not a fourth stability model. It asks about the small short-run residual change after last
year's corruption is considered. Its coefficient and within-
\Rtwo{} are not directly comparable to static M1--M3. Dynamic fixed-effects estimates may also exhibit
finite-panel bias \citep{nickell1981}; the relatively long country histories reduce, but do not remove,
that concern.

\subsection{Between-versus-within comparison}

To locate the raw positive relationship, the same single-predictor association is estimated as pooled
OLS, between-country OLS on country means, country fixed effects, and two-way fixed effects. Pooled
errors are clustered by country and between-country errors are heteroskedasticity-robust. The purpose
is not to select a winner by fit, but to show whether the sign comes from comparing different countries
or from changes within the same country.

\subsection{Democracy heterogeneity}

The hypothesis that institutional context moderates the association is tested by interacting populism
with electoral democracy:

\begin{equation}
  C_{it} = \beta_1 P_{it} + \beta_2 D_{it} + \beta_3(P_{it}D_{it})
  + \gamma\log GDPpc_{it} + \alpha_i + \lambda_t + \varepsilon_{it}.
\end{equation}

Rather than interpret $\beta_3$ alone, the report presents the model-implied slope
$\partial C/\partial P = \beta_1 + \beta_3D$ at the 10th, 50th, and 90th percentiles of democracy.
These are marginal effects from one interaction model, not separate regressions chosen after observing
the data.

\section{Temporal ordering}

\subsection{Lead-lag models}

RQ2 tests both possible directions. For populism to corruption, the two-lag model is

\begin{equation}
  C_{it} = \theta_1 P_{i,t-1} + \theta_2 P_{i,t-2}
  + \mathbf{X}_{it}'\boldsymbol{\gamma} + \alpha_i + \lambda_t + e_{it}.
  \label{eq:leadlag}
\end{equation}

The reverse model exchanges $C$ and $P$. Lead-lag estimates establish temporal order more clearly than
same-year association, but remain vulnerable to persistence. In the analysis panel, lag-1
autocorrelation is approximately 0.93 for populism and 0.99 for composite corruption.

\subsection{Granger-style tests}

The stricter model conditions on the dependent variable's own history:

\begin{equation}
  C_{it} = \sum_{\ell=1}^{2}\rho_{\ell} C_{i,t-\ell}
  + \sum_{\ell=1}^{2}\theta_{\ell} P_{i,t-\ell}
  + \mathbf{X}_{it}'\boldsymbol{\gamma} + \alpha_i + \lambda_t + e_{it}.
  \label{eq:granger}
\end{equation}

A joint Wald test evaluates $H_0:\theta_1=\theta_2=0$. Rejection means that lagged populism adds
predictive information beyond corruption's own past. It does not establish structural causality:
unobserved time-varying confounders, measurement error, and anticipation can still generate predictive
ordering. We therefore use ``Granger-style'' rather than causal language
\citep{granger1969}.

Lag windows from one to five years are compared with AIC and BIC only as diagnostics. Longer windows
drop additional early observations, so the information criteria do not always compare identical
samples. The two-lag model remains the pre-specified main analysis.

\subsection{Placebos, events, and changes}

Three designs probe the main directional result:

\begin{enumerate}
  \item \textbf{Placebo leads.} Future proposed-cause values are added to the current-outcome model.
  Significant future terms would warn of pre-trends, anticipation, or omitted common causes.
  \item \textbf{Event studies.} Populist entry is the first election-year increase in governing
  populism of at least 0.10 that reaches 0.50. A corruption shock is the first annual increase in the
  composite index of at least 0.05. Event-time coefficients are estimated relative to year $-1$, with
  joint pre- and post-event tests.
  \item \textbf{First differences.} Annual changes remove country levels through an alternative
  transformation. Directional models regress the current change in the outcome on the lagged change in
  the proposed cause and changes in controls.
\end{enumerate}

Because populism is forward-filled, the lead-lag model is also re-estimated on election years,
government-change years, and observations no more than five or ten years from the last election.

\section{Institutional decomposition}

For each corruption dimension $d$, the controlled static model is

\begin{equation}
  C^d_{it} = \beta_d P_{it} + \mathbf{X}_{it}'\boldsymbol{\gamma}_d
  + \alpha^d_i + \lambda^d_t + \varepsilon^d_{it}.
  \label{eq:decomposition}
\end{equation}

Raw coefficients cannot be ranked across outcomes with different latent scales. We therefore report
the within-standardised effect

\begin{equation}
  \widetilde{\beta}_d = \beta_d
  \frac{SD_{within}(P)}{SD_{within}(C^d)}.
\end{equation}

This quantity is the change in within-country standard deviations of the outcome associated with a
one-within-standard-deviation change in populism. Outcome-specific samples preserve power; a common
complete-case sample checks whether rankings reflect legislative missingness.

The primary decomposition contains four same-year M3 tests and eight Granger-style tests (four
dimensions in two directions). Benjamini--Hochberg adjustment controls the false discovery rate within
these pre-declared families \citep{benjamini1995}. The composite is a benchmark and is not counted as
an independent fifth corruption outcome.

\section{Predictive benchmark}

A gradient-boosted tree model provides a separate non-linear diagnostic. XGBoost uses governing-party
populism, log GDP per capita, and electoral democracy \citep{chen2016}. Entire countries are held out
with region-stratified group folds, preventing observations from the same country appearing in both
training and test sets. A nested inner loop tunes the hyperparameters without seeing outer-fold test
countries. SHAP values rank feature use in the fitted prediction model \citep{lundberg2017}.

This exercise is not an alternative causal estimator. Raw features retain between-country information
that fixed effects deliberately remove, so out-of-fold predictive \Rtwo{} and fixed-effects within-
\Rtwo{} answer different questions. The useful diagnostic is whether a flexible model finds populism
important after GDP and democracy are available.

\section{Decision rules and reproducibility}

The analysis follows four reporting rules. First, coefficients, intervals, sample sizes, and effect
magnitudes are reported together; statistical significance is never treated as substantive size.
Second, null findings are stated as absence of robust evidence, not proof of no effect. Third,
robustness checks evaluate stability and do not form a leaderboard from which the preferred result is
selected. Fourth, causal wording is reserved for designs that could justify it; this observational
panel cannot.

All preprocessing and analysis steps are executable from versioned Python modules. Tables and figures
are generated from committed processed data, and automated tests validate aggregation, orientation,
panel shifts, model construction, and multiplicity correction. Reproducibility commands are listed in
\Cref{app:reproducibility}.
